% Generated draft: requires mathematical and source review.
\documentclass[11pt]{article}
\usepackage[margin=25mm]{geometry}
\usepackage{fontspec}
\setmainfont{Latin Modern Roman}
\usepackage{amsmath,amssymb,mathtools,unicode-math}
\setmathfont{Latin Modern Math}
\usepackage{enumitem,needspace,xurl,hyperref}
\setlength{\parindent}{0pt}
\setlength{\parskip}{5pt}
\setlength{\emergencystretch}{3em}
\newcommand{\N}{\mathbb{N}}
\newcommand{\Z}{\mathbb{Z}}
\newcommand{\Q}{\mathbb{Q}}
\newcommand{\R}{\mathbb{R}}
\newcommand{\C}{\mathbb{C}}
\newcommand{\F}{\mathbb{F}}
\newcommand{\A}{\mathbb{A}}
\newcommand{\PP}{\mathbb{P}}
\newcommand{\E}{\mathbb{E}}
\newcommand{\eps}{\varepsilon}
\newcommand{\dd}{\,\mathrm{d}}
\begin{document}
\section*{UnsolvedMath ID: 2000007}
\textbf{Source problem code:} catalogue entry
\subsection{Definitions and mathematical statement}
In a standard elasticity formulation, a reference body occupies $\Omega\subset\mathbb R^d$ and a deformation is a map $y:\Omega\to\mathbb R^d$ with weak gradient $Dy$. For sufficiently regular orientation-preserving deformations, the Ciarlet--Ne\v{c}as noninterpenetration condition is written
\[
\int_{\Omega}\det Dy(x)\,dx\leq\mathcal L^d\big(y(\Omega)\big),
\qquad \det Dy>0\ \text{almost everywhere},
\]
where $\mathcal L^d$ is Lebesgue measure. Its connection with almost-everywhere injectivity requires appropriate regularity and change-of-variables hypotheses; it is not an assertion for arbitrary Sobolev representatives. Smooth self-contact conventionally permits distinct boundary portions of the deformed body to meet smoothly, without overlap of material interiors. It need not be compatible with demanding pointwise injectivity of the entire closed body.

\textbf{Problem.} Justify the Ciarlet--Ne\v{c}as minimization problem, or formulate and justify an appropriate modification, for situations involving smooth self-contact. Schematically the minimization concerns
\[
\inf_{y\in\mathcal A}\mathcal E(y),
\]
where $\mathcal E$ is the elastic total energy and $\mathcal A$ incorporates the intended boundary conditions, orientation preservation, and a noninterpenetration constraint permitting the specified contact. The supplied record does not define $\mathcal E$, $\mathcal A$, the contact regularity, or what result constitutes justification. It therefore supports this modeling-and-analysis target, not a uniquely specified existence or Euler--Lagrange theorem.

The supplied summary reports substantial progress through the Ciarlet--Ne\v{c}as formulation and weak-injectivity modifications, mentioning Ball's (INV) machinery and later work. It qualifies that no complete unification of existence and regularity under all natural assumptions is settled; no universal resolution is asserted here.
\subsection{Short English statement}
Give a mathematically justified elasticity minimization model that permits smooth contact between boundary portions while excluding interpenetration. The source leaves the energy, admissible class, and desired justification theorem unspecified.
\subsection{Sources}
\begin{itemize}
\item J. M. Ball, \emph{Some Open Problems in Elasticity} (2002), Problem 7, \url{https://people.maths.ox.ac.uk/ball/Articles%20in%20Conference%20Proceedings%20and%20Books/JMB%202002%20re%20Marsden%2060th.pdf}; PDF-text extraction, accessed 2026-07-29.
\item The problem names the Ciarlet--Ne\v{c}as formulation. The supplied summary additionally mentions Ball's (INV) machinery and subsequent works, but gives no separate paper, theorem, date, or precise version supporting its progress assessment.
\item Local registry \texttt{data/unsolvedmath/problems.json}, permanent ID $2000007$, source code catalogue entry, \emph{Some Open Problems in Elasticity --- Smooth self-contact}. The record labels this partially solved and PARTIAL-PROGRESS; its source-status and rights checks remain marked NEEDS\_REVIEW.
\end{itemize}
\end{document}
